
\documentclass{article} % For LaTeX2e
\usepackage{iclr2027_conference,times}

% Optional math commands from https://github.com/goodfeli/dlbook_notation.
\input{math_commands.tex}

\usepackage{hyperref}
\usepackage{url}
\usepackage{graphicx}
\usepackage{wrapfig}
\usepackage{float}
\usepackage{booktabs}
\usepackage{array}
\usepackage{capt-of}
\usepackage{listings}
\lstset{
  basicstyle=\ttfamily\footnotesize,
  breaklines=true,
  columns=fullflexible,
  frame=single,
  framerule=0.4pt,
  framesep=0.6em,
  xleftmargin=0.25em,
  xrightmargin=0.25em,
  aboveskip=0.75em,
  belowskip=0.75em,
  showstringspaces=false
}



\title{AI-H\textsuperscript{2}RAG: Asymmetric Indexing-based Hierarchical Hyper-RAG}

% Authors must not appear in the submitted version. They should be hidden
% as long as the \iclrfinalcopy macro remains commented out below.
% Non-anonymous submissions will be rejected without review.

\author{Anonymous authors\\
Paper under double-blind review}

% The \author macro works with any number of authors. There are two commands
% used to separate the names and addresses of multiple authors: \And and \AND.
%
% Using \And between authors leaves it to \LaTeX{} to determine where to break
% the lines. Using \AND forces a linebreak at that point. So, if \LaTeX{}
% puts 3 of 4 authors names on the first line, and the last on the second
% line, try using \AND instead of \And before the third author name.

\newcommand{\fix}{\marginpar{FIX}}
\newcommand{\new}{\marginpar{NEW}}

%\iclrfinalcopy % Uncomment for camera-ready version, but NOT for submission.
\begin{document}

\maketitle

\begin{abstract}
Hypergraph-based Retrieval-Augmented Generation(RAG) retrieves $n$-ary relational facts, but generally ignores the hierarchical organization of documents, through which information is conveyed across levels. This omission hinders the assembly of coherent evidence when answers require both global context and local facts. To address this limitation, we introduce a new framework called Asymmetric Indexing-based Hierarchical Hyper-RAG(AI-H\textsuperscript{2}RAG). It consists of a comprehensive pipeline, including: (i) hierarchical hypergraph construction, (ii) asymmetric indexing, and (iii) structure-aware retrieval and generation. It constructs a hierarchical hypergraph to connect global context with local facts, propagates information across its levels through asymmetric indexing, and aligns each query with the indexed hierarchical hypergraph to retrieve evidence in a coarse-to-fine manner for answer generation. Experiments across eight document question-answering benchmarks demonstrate that AI-H\textsuperscript{2}RAG outperforms standard, graph-based, and hypergraph-based RAG baselines in answer quality. Notably, our method achieves the strongest performance across questions requiring different evidence depths in multi-stage Mechanical question answering, while maintaining retrieval latency comparable to the strongest baselines.
\end{abstract}

\section{Introduction}
Retrieval-augmented generation (RAG) augments generation with external evidence retrieved at inference time \citep{Lewis2020RAG,Guu2020REALM}. By decoupling knowledge from model parameters, it supports source-grounded answers and updates to evolving corpora without retraining. These properties are particularly valuable in specialized domains, where answers depend on document-specific terminology, constraints, and time-varying evidence \citep{Yang2025KBSLLMKB}. Accordingly, RAG has emerged as a widely adopted paradigm in both academic research and industrial applications.

Current RAG methodologies broadly retrieve evidence from documents, graphs, or hypergraphs. Document-based RAG independently indexes text chunks and retrieves them through learned query--passage similarity \citep{Izacard2023Atlas}.Graph-based RAG organizes entities and relations into structured retrievers, connecting distributed evidence through entity paths or aggregating support within communities \citep{Jin2024LLMGraphs,Zhong2023KGConstruction}. Hypergraph-based RAG further represents high-order relations among multiple entities or facts \citep{Feng2019HGNN,Feng2026HyperRAG}. However, none of these paradigms explicitly captures the hierarchical document structure.

\begin{figure*}[!t]
\centering
\includegraphics[width=0.86\textwidth]{figs/Caschyper_overall_framework.pdf}
\caption{Overview of AI-H\textsuperscript{2}RAG. The document side constructs nested topic--chunk, chunk--sentence, and sentence--entity hypergraphs, then builds a context-aware index through asymmetric bidirectional fusion. The query side performs cascaded enhancement, topic-scoped chunk recall, sentence reranking, and multi-hop retrieval before generation.}
\label{fig:Caschyper_overall_framework}
\end{figure*}

Documents naturally organize evidence across several levels, including topics, chunks, sentences, and entities. Information is transmitted across these levels: topics define the scope of related chunks, chunks provide context for sentences, and sentences ground entities and their relations. Answering complex questions therefore requires jointly integrating evidence from the document hierarchy, rather than relying on isolated evidence at a single granularity. Hypergraphs are well suited to this setting because a hyperedge can directly connect multiple heterogeneous elements across levels, capturing the many-to-many dependencies among topics, chunks, sentences, and entities \citep{Chien2022AllSet}. Such a unified representation preserves both broad contextual scope and fine-grained factual specificity throughout the retrieval process.

Motivated by these insights, we design AI-H\textsuperscript{2}RAG, a hierarchical hypergraph RAG framework with asymmetric indexing. It is built upon three key steps. First, we propose a hierarchical hypergraph construction method, which builds three-layer nested hypergraphs to hierarchically represent the document structure. According to the document hierarchy, we represent the document as topic--chunk, chunk--sentence, and sentence--entity hypergraphs. This construction clearly supports information propagation across different levels of the document. Second, we introduce an asymmetric indexing mechanism. Within the constructed hierarchical hypergraph structure, it enables aggregated context flows downward, while original fine-grained facts flow upward to support higher-level semantics. Lastly, we develop a structure-aware retrieval and generation strategy that aligns each query with the indexed hierarchical hypergraph to retrieve evidence from coarse to fine granularity for answer generation.

To validate the effectiveness, We conduct experiments in multiple knowledge-intensive document QA benchmarks \citep{Guo2025LightRAG}, spanning physics, pathology, mix and five machine-tool domains, including Mechanical. Results demonstrate that AI-H\textsuperscript{2}RAG outperforms a range of previous RAG methods, and achieves the strongest performance across questions requiring different evidence depths in multi-stage Mechanical question answering, while maintaining retrieval latency comparable to the strongest baselines, demonstrating its great potential for practical, real-world applications.

Our contributions are the following:
\begin{itemize}
    \item We propose AI-H²RAG, a hierarchical hypergraph RAG framework that models documents as three nested hypergraph layers, bridging global context and local facts to answer multi-granularity evidence-based questions.
    \item We introduce an asymmetric indexing mechanism that propagates aggregated contextual representations flow downward from coarse to fine document levels, while original fine-grained factual representations flow upward to support coarse-level semantics.
    \item We develop a structure-aware retrieval and generation strategy that retrieves coarse-to-fine evidence by aligning queries with the hierarchical hypergraph for answer generation.
\end{itemize}

\section{Related Work}\label{sec:related-work}

\smallskip
\noindent\textbf{RAG}
Retrieval-augmented generation (RAG) couples a retriever with a language model to provide external evidence for knowledge-intensive tasks \citep{Asai2024SelfRAG}. Scaled retrieval and self-reflective variants extend this paradigm \citep{Borgeaud2022RETRO}. Recent surveys summarize the rapidly expanding RAG literature \citep{Gao2024RAGSurvey,Huang2026RAGSurvey}. Dense retrieval improves evidence access through learned query and passage representations. Query rewriting and hypothetical-document expansion improve underspecified retrieval inputs \citep{Ma2023QueryRewriting}. Question-answering benchmarks such as SQuAD and Natural Questions evaluate evidence-grounded answering \citep{Rajpurkar2016SQuAD,Kwiatkowski2019NaturalQuestions}. RAG evaluation also assesses retrieval grounding and generated responses \citep{Es2024RAGAS}. However, conventional RAG usually indexes passages independently, limiting its ability to coordinate document scope with local factual evidence \citep{Adjali2024MultiLevelRAG,Liu2024LostMiddle}.

\smallskip
\noindent\textbf{Graph-based RAG}
Graph-based RAG connects related evidence through entities, triples, and communities \citep{Pan2024LLMKG,Peng2025GraphRAGSurvey}. GraphRAG supports local and global query-focused summarization through entity graphs and community reports \citep{Edge2025GraphRAG}. LightRAG improves retrieval efficiency with dual-level graph retrieval \citep{Guo2025LightRAG}. Knowledge-graph methods also support relation-path reasoning for complex question answering \citep{Lan2023ComplexKBQA,Xu2024KGRAGCSQA}. Graph neural encoders aggregate neighborhood information through convolution or sampled message passing \citep{Kipf2017GCN,Hamilton2017GraphSAGE}. Attention-based graph encoders provide an alternative mechanism for weighting neighborhood information \citep{Velickovic2018GAT,Wu2021GNNSurvey}. Entity linking resolves surface mentions to canonical graph nodes \citep{Hoffart2011AIDA,Sevgili2022EntityLinkingSurvey}, while entity-aware contextual encoders provide learned entity representations \citep{Yamada2020LUKE}. Nevertheless, these methods primarily organize pairwise relations, entity neighborhoods, or coarse summaries, without explicitly constructing the document's hierarchical structure for multi-granularity retrieval needed for QA.

\smallskip
\noindent\textbf{Hypergraph-based RAG}
Hypergraphs generalize ordinary graphs by allowing each hyperedge to connect multiple vertices, thereby representing high-order dependencies beyond pairwise links \citep{Feng2019HGNN,Gao2023HGNNPlus}. Set-based hypergraph architectures provide a complementary formulation for aggregating variable-sized neighborhoods \citep{Chien2022AllSet}. Hyper-RAG extends this principle to retrieval-augmented generation through multi-entity hypergraph organization \citep{Feng2026HyperRAG}. Existing hypergraph RAG methods, however, typically construct hyperedges over entities or other single-level units. They do not explore hypergraphs as a means of explicitly representing the hierarchical structure of documents. Therefore, we propose AI-H\textsuperscript{2}RAG, which hierarchically represents document structure with hypergraphs to support multi-granularity evidence retrieval.

\section{Method}\label{sec:method}

In this section, we introduce the proposed method AI-H\textsuperscript{2}RAG, as shown in Figure~\ref{fig:Caschyper_overall_framework}, including hierarchical hypergraph construction, asymmetric indexing, and structure-aware retrieval and generation. 

\subsection{Hierarchical Hypergraph Construction}
\label{sec:hierarchical_hypergraph}

Documents organize evidence across topics, chunks, sentences, and entities. To represent document levels, we propose a hierarchical hypergraph construction method that includes multi-level hypergraph extraction, cross-level membership mapping, and raw representation initialization. The overall architecture is illustrated in Figure~\ref{fig:hypergraph_construction}.

\smallskip
\noindent\textbf{Multi-level Hypergraph Extraction.}
Given a document corpus $\mathcal{D}$, we decompose each document $d_n \in
\mathcal{D}$ into a set of chunks $\mathcal{C}$, sentences $\mathcal{S}$, and entities $\mathcal{U}$. Here, we use sentences to refer to sentence-level textual units, and employ the term entities broadly to encompass both conventional entities and salient semantic elements, including actions, conditions, events, and relations. The prompts for extracting sentences and entities are provided in Appendix~\ref{app:sentence_extraction_prompt} and Appendix~\ref{app:entity_extraction_prompt}. To represent these textual units and their relationships, a hypergraph is defined as $\mathcal{H}=(\mathcal{V},\mathcal{E})$, where $\mathcal{V}$ is the vertex set and $\mathcal{E}$ is the hyperedge set.

We first construct a fuzzy topic--chunk hypergraph. To obtain topic representations, we induce a set of latent topic anchors
$\mathcal{T}=\{t_k\}_{k=1}^{K}$ and a fuzzy topic--chunk membership matrix
$\mathbf{M}^{\mathrm{T}\to\mathrm{C}}=[m_{ik}]$ from the chunk set
$\mathcal{C}=\{c_i\}_{i=1}^{N}$, where $m_{ik}\in[0,1]$ denotes the membership
strength of chunk $c_i$ with respect to topic $t_k$. Unlike hard topic partitioning, this fuzzy assignment preserves cross-topic semantic associations by allowing a chunk to participate in multiple topic hyperedges. The detailed fuzzy-topic induction procedure is provided in Appendix~\ref{app:latent_topic_construction}.

Each topic $t_k$ serves as the semantic anchor of a topic-induced hyperedge
$e_k^{\mathrm{T}\to\mathrm{C}}=\{c_i\in\mathcal{C}\mid m_{ik}>0\}$.
Accordingly, with
$\mathcal{E}^{\mathrm{T}\to\mathrm{C}}
=
\{e_k^{\mathrm{T}\to\mathrm{C}}\}_{k=1}^{K}$,
we define the topic--chunk hypergraph as
\begin{equation}
\mathcal{H}^{\mathrm{T}\to\mathrm{C}}
=
\left(
\mathcal{C},
\mathcal{E}^{\mathrm{T}\to\mathrm{C}}
\right).
\end{equation}

We further construct two containment-based hypergraph layers:
\begin{equation}
\mathcal{H}^{\mathrm{C}\to\mathrm{S}}
=
\left(
\mathcal{S},
\mathcal{E}^{\mathrm{C}\to\mathrm{S}}
\right),
\qquad
\mathcal{H}^{\mathrm{S}\to\mathrm{U}}
=
\left(
\mathcal{U},
\mathcal{E}^{\mathrm{S}\to\mathrm{U}}
\right).
\label{eq:containment_hypergraphs}
\end{equation}
where each chunk $c_i$ is associated with a hyperedge $e_i^{\mathrm{C}\to\mathrm{S}}$ that connects its contained sentences, and each sentence $s_j$ is associated with a hyperedge $e_j^{\mathrm{S}\to\mathrm{U}}$ that connects its extracted entities. Consequently, the three nested hypergraph layers represent four document levels: topics, chunks, sentences, and entities.

\smallskip
\noindent\textbf{Cross-level Membership Mapping.}
We explicitly preserve the membership relations between adjacent document
levels using three incidence mappings:
$\mathbf{M}^{\mathrm{T}\to\mathrm{C}}=[m_{ik}]\in[0,1]^{N\times K}$,
$\mathbf{M}^{\mathrm{C}\to\mathrm{S}}$, and
$\mathbf{M}^{\mathrm{S}\to\mathrm{U}}$.
The topic--chunk mapping retains fuzzy membership weights, whereas the
chunk--sentence mapping encodes explicit containment. The sentence--entity
mapping further retains the extraction-time importance of each entity:
\begin{equation}
m_{ij}^{\mathrm{C}\to\mathrm{S}}
=
\mathbb{I}[s_j \subset c_i],
\qquad
m_{jm}^{\mathrm{S}\to\mathrm{U}}
=
\rho_m\,
\mathbb{I}[u_m \subset s_j],
\label{eq:containment_membership_mapping}
\end{equation}
where $\rho_m\in[0,1]$ is the importance score assigned to entity $u_m$ during semantic-unit extraction. These mappings allow retrieval to move structurally in both directions. A fine-grained entity can be traced back to its enclosing sentence, chunk, and topic, while a coarse topic can look for local facts that support it

\smallskip
\noindent\textbf{Raw Representation Initialization.}
Before cross-level information propagation, we initialize an $\ell_2$-normalized raw representation for every document unit. Let $\operatorname{Enc}_{\mathrm{bge}}(\cdot)$ denote the BGE-M3 text encoder and $\operatorname{norm}(\cdot)$ denote $\ell_2$ normalization. For chunks, sentences, and entities, we obtain
\begin{equation}
\mathbf{h}_{c_i}^{(0)}
=
\operatorname{norm}\!\left(
\operatorname{Enc}_{\mathrm{bge}}(c_i)
\right),
\hspace{1.2em}
\mathbf{h}_{s_j}^{(0)}
=
\operatorname{norm}\!\left(
\operatorname{Enc}_{\mathrm{bge}}(s_j)
\right),
\hspace{1.2em}
\mathbf{h}_{u_m}^{(0)}
=
\operatorname{norm}\!\left(
\operatorname{Enc}_{\mathrm{bge}}(u_m)
\right).
\label{eq:raw_representations}
\end{equation}

The initial representation of topic $t_k$ is obtained by aggregating its member chunks according to their fuzzy memberships:
\begin{equation}
\mathbf{h}_{t_k}^{(0)}
=
\operatorname{norm}\!\left(
\frac{
\sum_{c_i\in e_k^{\mathrm{T}\to\mathrm{C}}}
m_{ik}\mathbf{h}_{c_i}^{(0)}
}{
\sum_{c_i\in e_k^{\mathrm{T}\to\mathrm{C}}}m_{ik}
}
\right).
\label{eq:topic_raw_representation}
\end{equation}

These representations preserve the original information available at each level.

\subsection{Asymmetric Indexing}

After constructing hierarchical hypergraph construction, We introduce an asymmetric indexing mechanism that .

\subsubsection{Local attention fusion}

Each semantic unit has a normalized raw embedding. To combine a target embedding $\mathbf{q}$ with a local feature set $K=\{\mathbf{k}_1,\ldots,\mathbf{k}_m\}$, we use the parameter-free fusion operator $\mathbf{F}_{attn}$. It assigns temperature-scaled similarity weights,
\begin{equation}
\alpha_i=\frac{\exp(\mathbf{q}^{\top}\mathbf{k}_i/\tau)}{\sum_{j=1}^{m}\exp(\mathbf{q}^{\top}\mathbf{k}_j/\tau)},
\qquad
\mathbf{F}_{attn}(\mathbf{q},K)=\operatorname{norm}\!\left(\sum_{i=1}^{m}\alpha_i\mathbf{k}_i\right),
\end{equation}
where $\tau$ controls the attention sharpness and $\operatorname{norm}$ denotes $L_2$ normalization. The operator retains the similarity-based form of dot-product attention \citep{Vaswani2017Attention}, but omits trainable projections. It uses the target unit's raw embedding as a reference vector for weighting local context, thereby selecting relevant information while preserving the unit's original semantic anchor.

% Floating wraps defer the graphic when necessary while prose keeps flowing.
\subsubsection{Index construction}

\begin{wrapfigure}{L}{0.50\textwidth}
\centering
\includegraphics[width=0.98\linewidth]{figs/asymmetric_bidirectional_cascaded_index.pdf}
\caption{Asymmetric bidirectional index construction. Final parent representations flow downward as context, while only raw child representations flow upward as factual support. Each fusion operation remains anchored by the target's raw representation.}
\label{fig:asymmetric_index}
\end{wrapfigure}

Figure~\ref{fig:asymmetric_index} shows the index construction. The cascade proceeds from chunks to sentences to entities. At each level, raw child representations provide bottom-up factual support, whereas final parent representations provide top-down contextual guidance. This asymmetry prevents contextual information from being repeatedly fed back as factual evidence in subsequent cascade updates.

According to Appendix~\ref{app:latent_topic_construction}, let $\mathcal{K}(c_i)=\{k\mid c_i\in e_k^{\mathrm{T}\to\mathrm{C}}\}$ denote the set of topic indices associated with chunk $c_i$. For each $k\in\mathcal{K}(c_i)$, $m_{ik}$ denotes the fuzzy membership of $c_i$ in topic $k$, and $\mathbf{t}_k$ denotes the corresponding topic-anchor representation. We first form the membership-weighted topic context
\begin{equation}
\mathbf{v}_{\mathrm{topic}}=
\operatorname{norm}\!\left(
\frac{\sum_{k\in\mathcal{K}(c_i)}m_{ik}\mathbf{t}_k}
{\sum_{k\in\mathcal{K}(c_i)}m_{ik}}
\right).
\end{equation}
Let $\mathbf{h}_{s\_agg}^{(raw)}$ be the mean raw representation of sentences in $c$. The final chunk representation is
\begin{equation}
\mathbf{h}_{c}^{(final)}=
\mathbf{F}_{attn}\!\bigl(\mathbf{h}_{c}^{(raw)},
\{\mathbf{v}_{\mathrm{topic}},\mathbf{h}_{s\_agg}^{(raw)},\mathbf{h}_{c}^{(raw)}\}\bigr).
\end{equation}

A sentence can occur in overlapping chunks. We therefore aggregate the final representations of all parent chunks $\mathcal{P}(s)$,
\begin{equation}
\mathbf{h}_{c\_agg}^{(final)}(s)=
\frac{1}{|\mathcal{P}(s)|}\sum_{c\in\mathcal{P}(s)}\mathbf{h}_{c}^{(final)}.
\end{equation}
Let $\mathbf{h}_{e\_agg}^{(raw)}(s)$ be the importance-weighted mean of the raw entities in $s$. The sentence update combines parent context, entity support, and its raw anchor:
\begin{equation}
\mathbf{h}_{s}^{(final)}=
\mathbf{F}_{attn}\!\left(
\mathbf{h}_{s}^{(raw)},
\{\mathbf{h}_{c\_agg}^{(final)}(s),\mathbf{h}_{e\_agg}^{(raw)}(s),\mathbf{h}_{s}^{(raw)}\}
\right).
\end{equation}
Finally, an entity has no lower-level children. It receives only its final parent-sentence context while retaining its raw representation:
\begin{equation}
\mathbf{h}_{e}^{(final)}=
\mathbf{F}_{attn}\!\left(
\mathbf{h}_{e}^{(raw)},
\{\mathbf{h}_{s}^{(final)},\mathbf{h}_{e}^{(raw)}\}
\right).
\end{equation}
Because the parent sentence already incorporates chunk- and topic-level context, this update can assign distinct representations to the same surface entity in different document contexts. Appendix~\ref{app:same_entity_different_sense_disambiguation} provides a contextual entity-disambiguation diagnostic for this capability.
Overall, the completed index couples broad context with local factual signals without requiring supervised graph training.

\subsection{Query-side Cascaded Multi-granularity Enhancement}

\begin{wrapfigure}{R}{0.50\textwidth}
\centering
\includegraphics[width=0.98\linewidth]{figs/Multi-hop retrieval and generation.pdf}
\caption{Multi-granularity retrieval and generation. Topic-guided chunk recall restricts the search space, sentence reranking selects direct evidence, and shared-entity expansion retrieves multi-hop support for generation.}
\label{fig:multihop_retrieval_generation}
\end{wrapfigure}

The document index is structurally rich, whereas an input query is often short and underspecified. We therefore expand the query into \(\mathbf{q}^{(exp)}\), extract its entity set \(\mathcal{E}_q\), and construct query representations that correspond to the document levels. Query expansion follows the motivation of hypothetical-document methods, which reduce the semantic mismatch between query and document representations \citep{Gao2023HyDE}. Figure~\ref{fig:multihop_retrieval_generation} provides an overview of the resulting retrieval and generation process.

The expanded query first activates relevant topic anchors. Its score for topic $k$ is $r_{qk}={\mathbf{q}^{(exp)}}^{\top}\mathbf{t}_k$. Given $r_q^{\max}=\max_k r_{qk}$, the active set is
\begin{equation}
\mathcal{A}_q=\{k\mid r_{qk}\geq\lambda r_q^{\max}\}.
\end{equation}
Their score-weighted mean defines the topic context $\mathbf{q}_{topic}$, and its associated topic--chunk hyperedges restrict the candidate chunks to a local subgraph. The raw expanded query remains the shared anchor for the chunk- and sentence-level representations:
\begin{equation}
\mathbf{q}_{c}=\mathbf{F}_{attn}\!\left(
\mathbf{q}^{(exp)},\{\mathbf{q}_{topic},\mathbf{q}^{(exp)}\}\right),
\end{equation}
\begin{equation}
\mathbf{q}_{s}=\mathbf{F}_{attn}\!\left(
\mathbf{q}^{(exp)},\{\mathbf{q}_{c},\mathbf{q}_{e\_agg},\mathbf{q}^{(exp)}\}\right),
\end{equation}
where $\mathbf{q}_{e\_agg}$ is the importance-weighted mean of raw query-entity representations. For each query entity $e_q\in\mathcal{E}_q$, we retain its raw representation as the fusion anchor and condition it on the final sentence-level query representation:
\begin{equation}
\mathbf{h}_{e_q}^{(final)}=\mathbf{F}_{attn}\!\left(
\mathbf{h}_{e_q}^{(raw)},\{\mathbf{q}_{s},\mathbf{h}_{e_q}^{(raw)}\}\right).
\end{equation}
The final query-entity set is
\begin{equation}
\mathbf{E}_q^{(final)}=\{\mathbf{h}_{e_q}^{(final)}\mid e_q\in\mathcal{E}_q\}.
\end{equation}

We first rank candidate chunks by
\begin{equation}
r_c={\mathbf{q}_{c}}^{\top}\mathbf{h}_{c}^{(final)},
\end{equation}
and retain the top-$K$ chunks. Within this scope, we rerank candidate sentences. Each sentence receives a contextual chunk score $s_c={\mathbf{q}^{(exp)}}^{\top}\mathbf{h}_{c}^{(final)}$ and a direct sentence score $s_s={\mathbf{q}^{(exp)}}^{\top}\mathbf{h}_{s}^{(final)}$. We use $\mathbf{q}^{(exp)}$ for these scores because it remains the stable lexical-semantic anchor before contextual query fusion. Appendix~\ref{app:query_scoring_comparison} compares this choice with cascaded query representations.

Let $\mathcal{E}_s$ denote the entity set extracted from a candidate sentence and $\mathbf{E}_s^{(final)}$ its final representations. The entity score retains fine-grained alignments through MaxSim \citep{Santhanam2022ColBERTv2}:
\begin{equation}
s_e=\frac{1}{|\mathbf{E}_q^{(final)}|}
\sum_{e_q\in\mathcal{E}_q}
\max_{e_d\in\mathcal{E}_s}
{\mathbf{h}_{e_q}^{(final)}}^{\top}\mathbf{h}_{e_d}^{(final)}.
\end{equation}
We obtain adaptive granularity weights from $[s_c,s_s,s_e]$ using a temperature-scaled Softmax. The reranking score is
% Retain the template's natural display spacing without page-bottom stretching.
\begingroup
\begin{equation}
\mathcal{S}^{(1)}=w_cs_c+w_ss_s+w_es_e.
\end{equation}
\endgroup
The top-$p$ sentences form the initial evidence set.

To retrieve distributed support, we alternate between selected sentences and their entities. Each entity is used to look up unvisited sentences sharing that entity; their entities then seed the next expansion step. This procedure recovers evidence chains across physical chunks, as required in multi-hop question answering \citep{Yang2018HotpotQA}. The final generator performs dual-track contrastive generation using two complementary evidence tracks: top-ranked chunks establish topical scope and background, while reranked and multi-hop sentences provide fine-grained support. Multi-hop candidates are retained only when they concern the query entities and remain consistent with the direct sentence and chunk evidence. The answer is consequently grounded in mutually supported contextual and factual evidence.

\section{Experiments}\label{sec:experiments}
\subsection{Experimental Setup}

We evaluate eight document QA datasets spanning five machine-tool and three general domains. The Mechanical evaluation uses three evidence-depth stages: Stage~1 requires direct evidence, Stage~2 requires one intermediate connection, and Stage~3 requires two connected evidence links. Baselines include an LLM without retrieval, dense RAG \citep{Karpukhin2020DPR}, GraphRAG \citep{Edge2025GraphRAG}, LightRAG \citep{Guo2025LightRAG}, HippoRAG~2 \citep{Gutierrez2025HippoRAG2}, and Hyper-RAG \citep{Feng2026HyperRAG}. Within each comparison and evaluation run, all methods share the documents, questions, reference evidence, generator, and evaluation protocol. GPT-4o-mini generates questions and judges answers against reference evidence on six criteria scored from 0 to 100. We conduct five independent evaluation runs. In each run, we randomly generate a distinct 50-question evaluation set for each dataset. We report the mean $\pm$ standard deviation across the five run-level mean scores; the average score is the mean across the six criteria. Appendix~\ref{app:experimental_setup} provides dataset statistics, preprocessing, and the detailed evaluation protocol; Appendix~\ref{app:llm_judge_scoring_prompt} reports the judging prompt. Appendices~\ref{app:expert_blind_evaluation} and~\ref{app:key_evidence_chain_retrieval} report the expert blind evaluation and evidence-chain retrieval analysis, respectively.

\subsection{Overall Performance Across RAG Methods}

At Stage~1 on the mechanical-domain datasets, Figure~\ref{fig:mechanical_llm_overall} compares AI-H\textsuperscript{2}RAG with other RAG methods across five LLM backbones: GLM-4.7, Kimi-K2.5, Qwen3.5-Flash, GPT-4o-mini, and DeepSeek-V3.2. AI-H\textsuperscript{2}RAG ranks first under every backbone, with average scores of 77.2, 85.7, 88.3, 86.5, and 86.7. Its relative improvement over the strongest baseline is 1.4--3.4\% across the five backbones. Absolute scores vary with the generator, from 77.2 under GLM-4.7 to 88.3 under Qwen3.5-Flash, but the ordering remains unchanged. Thus, the retrieval gain persists across generators with different answer capabilities.

Mechanical provides a controlled comparison across evidence depths: Stage~1 uses direct evidence, Stage~2 one intermediate connection, and Stage~3 two connected links (Appendix Table~\ref{tab:mechanical_question_examples}). With the generator fixed, AI-H\textsuperscript{2}RAG scores 86.5, 88.9, and 88.1 on Stages~1--3, yielding relative improvements of 3.3\%, 3.0\%, and 3.4\%, respectively, over the strongest baseline, Hyper-RAG (Figure~\ref{fig:mechanical_stage_overall}). The gains therefore persist across direct and multi-hop questions.

% Keep the two overall-performance figures in numerical and narrative order.
\clearpage
\begin{figure}[H]
\centering
\includegraphics[width=0.94\textwidth]{figs/summary_1_stage_mechanical_grouped_bars.pdf}
\caption{Six evaluation-dimension scores and average scores of different RAG methods under different LLM backbones. Error bars denote the standard deviation across five run-level mean scores.}
\label{fig:mechanical_llm_overall}
\end{figure}

\begin{figure}[H]
\centering
\includegraphics[width=0.98\textwidth]{figs/gpt4o_mini_mechanical_stage_difficulty.pdf}
\caption{Six evaluation-dimension scores and average scores of different RAG methods across Mechanical question stages. Each stage contains 50 questions per evaluation run; error bars denote the standard deviation across five run-level mean scores.}
\label{fig:mechanical_stage_overall}
\end{figure}

\subsection{Cross-domain Generalization}

% Match the wrapped lines to the actual graphic and caption height.
\begin{wrapfigure}{R}{0.52\textwidth}
\centering
\includegraphics[width=0.98\linewidth]{figs/gpt4o_mini_dataset_average_line.pdf}
\caption{Cross-domain average scores under the same GPT-4o-mini backbone.}
\label{fig:cross_domain_average_line}
\end{wrapfigure}

We fix the generator and compare LightRAG, Hyper-RAG, and AI-H\textsuperscript{2}RAG across all eight corpora. Figures~\ref{fig:cross_domain_average_line} and~\ref{fig:cross_domain_overall} show that AI-H\textsuperscript{2}RAG achieves the highest average score in every domain.

Its relative improvements over the strongest baseline are 3.3\%, 1.7\%, 6.6\%, 7.6\%, and 3.5\% on the five machine-tool corpora, and 3.5\%, 3.2\%, and 3.0\% on Mix, Pathology, and Physics. The radar profiles show gains across all six criteria rather than a gain driven by a single \mbox{dimension}. All methods score lower on Mix, which contains more heterogeneous terminology and evidence patterns, yet AI-H\textsuperscript{2}RAG still ranks first. These results support transfer across both specialized technical and general-domain corpora.

\par\WFclear
\begin{figure}[H]
\centering
\includegraphics[width=0.95\textwidth]{figs/gpt4o_mini_dataset_method_radar.pdf}
\caption{Six-dimensional cross-domain comparison of LightRAG, Hyper-RAG, and AI-H\textsuperscript{2}RAG across eight datasets under the same GPT-4o-mini backbone.}
\label{fig:cross_domain_overall}
\end{figure}

\subsection{Ablation Study}

We ablate the three proposed designs on 50 Stage~1 Mechanical questions per evaluation run: cascaded multi-granularity hypergraph construction, asymmetric bidirectional indexing with local attention fusion, and query-side cascaded enhancement. For the indexing innovation, we use two complementary replacements: raw-parent propagation replaces contextual propagation from final parent representations to child representations, and mean fusion replaces local attention fusion. Each ablation replaces one design while keeping the remaining pipeline and evaluation protocol fixed.
\begin{table*}[t]
\centering
\caption{Ablations on 50 Stage~1 Mechanical questions per evaluation run using GPT-4o-mini. Metrics are mean $\pm$ SD across five run-level means; $\Delta$ is the relative mean-score change from the full model.}
\label{tab:ablation_results}
\setlength{\tabcolsep}{3pt}
\scriptsize
{\renewcommand{\arraystretch}{1.10}
\resizebox{\textwidth}{!}{%
\begin{tabular}{@{}>{\scriptsize\raggedright\arraybackslash}p{0.34\textwidth}rrrrrrrr@{}}
\toprule
Ablated design and fallback & \multicolumn{1}{c}{Comp.} & \multicolumn{1}{c}{Div.} & \multicolumn{1}{c}{Emp.} & \multicolumn{1}{c}{Log.} & \multicolumn{1}{c}{Read.} & \multicolumn{1}{c}{Rel.} & \multicolumn{1}{c}{Avg.} & \multicolumn{1}{c}{$\Delta$ (\%)} \\
\midrule
None (full AI-H\textsuperscript{2}RAG) & $91.04 \pm 0.83$ & $81.94 \pm 0.85$ & $79.76 \pm 0.91$ & $86.00 \pm 0.37$ & $85.36 \pm 0.27$ & $94.80 \pm 0.32$ & \textbf{$86.48 \pm 0.38$} & 0.0\% \\
Cascaded multi-granularity hypergraph (raw-unit index) & $83.80 \pm 1.21$ & $71.90 \pm 1.05$ & $75.70 \pm 0.73$ & $83.60 \pm 0.44$ & $84.10 \pm 0.31$ & $86.61 \pm 0.29$ & $80.95 \pm 0.42$ & -6.4\% \\
Asymmetric final-parent propagation (raw-parent propagation) & $86.00 \pm 1.24$ & $71.70 \pm 1.37$ & $75.70 \pm 0.84$ & $85.50 \pm 0.21$ & $83.20 \pm 0.44$ & $92.40 \pm 0.23$ & $82.42 \pm 0.58$ & -4.7\% \\
Local attention fusion (mean fusion) & $84.80 \pm 1.03$ & $75.70 \pm 1.44$ & $76.30 \pm 0.65$ & $84.60 \pm 0.34$ & $83.42 \pm 0.38$ & $91.54 \pm 0.21$ & $82.73 \pm 0.38$ & -4.3\% \\
Query-side cascaded enhancement (raw query representations) & $84.20 \pm 0.97$ & $69.40 \pm 1.39$ & $75.50 \pm 0.78$ & $83.90 \pm 0.33$ & $84.72 \pm 0.37$ & $89.20 \pm 0.25$ & $81.15 \pm 0.54$ & -6.2\% \\
\bottomrule
\end{tabular}%
}
}
\end{table*}

Table~\ref{tab:ablation_results} shows a decrease for every replacement. Replacing the cascaded hypergraph with a raw-unit index reduces the average score by 6.4\%, and replacing query-side enhancement with raw query representations reduces it by 6.2\%. Raw-parent propagation and mean fusion reduce the score by 4.7\% and 4.3\%, respectively. The results indicate that the benefit is not explained by retaining multiple text units alone: cascaded organization, contextual indexing, and query-side alignment each contribute.

\subsection{Retrieval Efficiency}
\label{sec:retrieval_efficiency}

We measure online retrieval and context-construction time on 50 Stage~1 Mechanical questions per evaluation run. Offline index construction and final answer generation are excluded. Each point in Figure~\ref{fig:retrieval_efficiency} combines the mean latency of a retrieval method with its average answer-quality score.

\begin{figure}[H]
\centering
\includegraphics[width=0.58\textwidth]{figs/retrieval_efficiency_mechanical.pdf}
\caption{Average retrieval efficiency on 50 first-stage Mechanical questions per evaluation run. Each point reports five run-level mean scores: average retrieval time and the corresponding average answer-quality score of a retrieval method.}
\label{fig:retrieval_efficiency}
\end{figure}

RAG is fastest because it directly matches text chunks. LightRAG is also relatively efficient, whereas Hyper-RAG and AI-H\textsuperscript{2}RAG spend more time on structured retrieval. AI-H\textsuperscript{2}RAG achieves the highest answer-quality score while retaining latency comparable to Hyper-RAG. Its additional topic-guided recall, sentence reranking, and multi-hop expansion therefore improve evidence organization without a large latency increase over the strongest structured-retrieval baseline.

\section{Conclusion}\label{sec:conclusion}

AI-H\textsuperscript{2}RAG organizes document evidence as nested topic--chunk, chunk--sentence, and sentence--entity hypergraphs, constructs an asymmetric bidirectional index, and applies cascaded multi-granularity enhancement to queries. Across eight document QA datasets, five LLM backbones, and three Mechanical question stages, it achieves the highest average performance. AI-H\textsuperscript{2}RAG achieves the highest answer-quality score while retaining retrieval latency comparable to Hyper-RAG. These results support the effectiveness of cascaded multi-granularity retrieval for answer quality and evidence coverage.

\clearpage
\section*{AI Use Statement}

In this work, we used GPT-based generative AI tools to assist with code development and literature organization. We also used GPT-4o-mini to generate synthetic evaluation questions from sampled corpus chunks and to automatically evaluate generated answers against the corresponding reference evidence using a predefined six-criterion rubric. These uses were part of the experimental protocol described in the Experimental Setup and Appendix. All AI-assisted code, literature materials, and experimental outputs were reviewed by the authors. The authors take responsibility for the final manuscript, experimental design, reported results, and all claims in this submission.



% Keep the bibliography's standard 4 pt item spacing fixed so a short final
% reference page does not stretch the gaps between entries.
\begingroup
\setlength{\bibsep}{4pt}
\bibliography{iclr2027_conference}
\bibliographystyle{iclr2027_conference}
\endgroup

\appendix
\raggedbottom
\numberwithin{equation}{section}
\numberwithin{figure}{section}
\numberwithin{table}{section}

\section{Cascaded Multi-Granularity Hypergraph Construction}
\label{app:cascaded_hypergraph_construction}

Figure~\ref{fig:hypergraph_construction} illustrates the three nested sub-hypergraphs used in the main-text construction. The topic--chunk level links semantically related document regions, while the chunk--sentence and sentence--entity levels retain local evidence relationships.

\begin{figure}[H]
\centering
\includegraphics[width=0.60\linewidth]{figs/cascaded_multigranularity_hypergraph_construction.pdf}
\caption{Cascaded multi-granularity hypergraph construction. Topics connect chunks, chunks connect sentences, and sentences connect entities.}
\label{fig:hypergraph_construction}
\end{figure}

\section{Latent Topic Construction from Chunk Similarity}
\label{app:latent_topic_construction}

This appendix details how the latent topic anchors $\mathcal{T}=\{t_k\}_{k=1}^{K}$ and the fuzzy topic--chunk membership matrix $\mathbf{M}^{\mathrm{T}\to\mathrm{C}}=[m_{ik}]$ are induced from the chunk set $\mathcal{C}=\{c_i\}_{i=1}^{N}$. As defined in Equation~\ref{eq:raw_representations}, $\mathbf{h}_{c_i}^{(0)}\in\mathbb{R}^{d}$ denotes the $\ell_2$-normalized raw representation of chunk $c_i$. The pairwise semantic similarity between chunks is therefore computed by
\begin{equation}
s_{ij}={\mathbf{h}_{c_i}^{(0)}}^{\top}\mathbf{h}_{c_j}^{(0)},
\quad i\neq j.
\end{equation}
To retain only informative semantic connections, we compute a dynamic threshold from the upper-triangular similarity values:
\begin{equation}
\tau=\operatorname{mean}\left(\{s_{ij}\mid i<j\}\right)
+0.5\operatorname{std}\left(\{s_{ij}\mid i<j\}\right).
\end{equation}
We then construct a weighted undirected chunk-similarity graph $\mathcal{G}=(\mathcal{C},\mathcal{R})$, where
\begin{equation}
(c_i,c_j)\in\mathcal{R}\ \Longleftrightarrow\ s_{ij}>\tau,
\qquad w_{ij}=s_{ij}.
\end{equation}
Leiden community detection is applied to $\mathcal{G}$ with edge weights $w_{ij}$ and resolution parameter $1.0$, producing a partition $\{\mathcal{C}_k\}_{k=1}^{K}$ of the chunks \citep{Traag2019Leiden}. Each community $\mathcal{C}_k$ induces one latent topic $t_k$; together, these topics form $\mathcal{T}$. The representation of topic anchor $t_k$ is the normalized centroid of its member chunk representations:
\begin{equation}
\boldsymbol{\mu}_k=\frac{1}{|\mathcal{C}_k|}\sum_{c_i\in\mathcal{C}_k}\mathbf{h}_{c_i}^{(0)},
\qquad
\mathbf{t}_k=\frac{\boldsymbol{\mu}_k}{\|\boldsymbol{\mu}_k\|_2}.
\end{equation}

Although Leiden provides a primary topic assignment, we make the topic--chunk memberships fuzzy to preserve cross-topic associations. Given a fuzzy-membership threshold $\eta$, we define
\begin{equation}
 m_{ik}=
\begin{cases}
1, & c_i\in\mathcal{C}_k,\\
{\mathbf{h}_{c_i}^{(0)}}^{\top}\mathbf{t}_k, & c_i\notin\mathcal{C}_k\ \text{and}\ {\mathbf{h}_{c_i}^{(0)}}^{\top}\mathbf{t}_k>\eta,\\
0, & \text{otherwise}.
\end{cases}
\end{equation}
Collecting these memberships yields $\mathbf{M}^{\mathrm{T}\to\mathrm{C}}=[m_{ik}]\in[0,1]^{N\times K}$. Thus, each chunk is assigned membership one in its primary Leiden community. It can additionally belong to other topic hyperedges when its similarity to their topic anchors exceeds the threshold $\eta$. The topic-induced hyperedge anchored by $t_k$ is
\begin{equation}
e_k^{\mathrm{T}\to\mathrm{C}}=\left\{c_i\in\mathcal{C}\mid m_{ik}>0\right\}.
\end{equation}
The resulting hyperedge set $\mathcal{E}^{\mathrm{T}\to\mathrm{C}}=\{e_k^{\mathrm{T}\to\mathrm{C}}\}_{k=1}^{K}$ therefore defines the topic--chunk hypergraph $\mathcal{H}^{\mathrm{T}\to\mathrm{C}}=(\mathcal{C},\mathcal{E}^{\mathrm{T}\to\mathrm{C}})$. A chunk may participate in multiple topic hyperedges when it has sufficiently high similarity to multiple topic anchors.

\section{Sentence-level Semantic-Unit Extraction Prompt}
\label{app:sentence_extraction_prompt}

The following prompt is used to extract the sentence-level semantic units that serve as vertices in the chunk--sentence hypergraph.
\begin{lstlisting}
SENTENCE_EXTRACTION_PROMPT = """
Extract all self-contained semantic units from the following paragraph. A semantic unit may contain one or more consecutive sentences when they together express one complete idea. Do not split sentences that depend on each other for essential context, and preserve the original wording and order.
Paragraph: {paragraph}
Return JSON only: {{"sentences": [{"text": "..."}, ...]}}
"""
\end{lstlisting}

\section{Entity Extraction Prompt}
\label{app:entity_extraction_prompt}

The following prompt is used to extract the entity-level semantic units that serve as vertices in the sentence--entity hypergraph.
\begin{lstlisting}
ENTITY_EXTRACTION_PROMPT = """
Extract the important information units from the sentence below. Identify both concrete entities and abstract semantic units that are useful for retrieval.

Sentence: {sentence}

Extract these types of information when they are explicitly present:
1. Concrete entities: people, organizations, locations, documents, products, methods, events, dates, and numerical values.
2. Abstract semantic units: central actions, goals, outcomes, causes, conditions, mechanisms, evidence, and query intent.

Importance scoring:
- 1.0: the main subject, core action, primary query intent, or indispensable result.
- 0.7-0.9: essential objects, key conditions, named entities, methods, evidence, or constraints.
- 0.4-0.6: useful modifiers, supporting details, and contextual attributes.
- 0.1-0.3: only include minor background details when they materially clarify the sentence.

Rules:
- Extract only text spans that occur in the sentence; do not invent entities or facts.
- Use concise, consistent type labels such as Person, Organization, Location, Document, Product, Event, Action, Goal, Method, Evidence, Condition, or QueryIntent.
- Assign a short description of each unit's role in the sentence.

Example 1 :
Input: "The Intervale supports new farmers through its Farms Program, which leases land, equipment, and storage facilities to commercial farmers. It also provides technical assistance and mentorship from experienced farmers."
Output:
{{"entities": [
  {{"name": "support new farmers", "type": "Action", "importance": 1.0, "description": "The central support action"}},
  {{"name": "Intervale", "type": "Organization", "importance": 0.9, "description": "The organization providing support"}},
  {{"name": "Farms Program", "type": "Program", "importance": 0.8, "description": "The support program"}},
  {{"name": "technical assistance and mentorship", "type": "Method", "importance": 0.8, "description": "Support mechanisms for new farmers"}}
]}}

Example 2 :
Input: "CytoSorb has received FDA Breakthrough Designation for the removal of ticagrelor in a cardiopulmonary bypass circuit during emergent and urgent cardiothoracic surgery. It also received Emergency Use Authorization for use in critically ill patients with COVID-19 infection."
Output:
{{"entities": [
  {{"name": "CytoSorb", "type": "Product", "importance": 1.0, "description": "The product being evaluated"}},
  {{"name": "FDA Breakthrough Designation", "type": "RegulatoryApproval", "importance": 1.0, "description": "A key regulatory milestone"}},
  {{"name": "Emergency Use Authorization", "type": "RegulatoryApproval", "importance": 1.0, "description": "An emergency regulatory authorization"}},
  {{"name": "ticagrelor", "type": "Drug", "importance": 0.8, "description": "The substance targeted for removal"}},
  {{"name": "COVID-19 infection", "type": "Condition", "importance": 0.8, "description": "The condition of the authorized patient group"}}
]}}

Return exactly one JSON object and no additional text:
{{"entities": [...]}}
"""
\end{lstlisting}

\section{Diagnostic for Same-Entity Different-Sense Disambiguation Across Document Contexts}
\label{app:same_entity_different_sense_disambiguation}

To provide a targeted diagnostic for same-entity different-sense disambiguation across document contexts, we constructed a balanced set of 150 entity pairs from the fixed multi-domain corpus, comprising 75 same-sense pairs and 75 different-sense pairs. Each pair contains two occurrences of the same entity in different document contexts. Each occurrence was used only as a query-time context over the fixed corpus index. We treated different-sense pairs as positive instances and used cosine distance between the two contextualized representations as the disambiguation score. We report AUROC with 95\% confidence intervals obtained from 2,000 stratified bootstrap samples.

We compare a Raw condition, which uses only the static embedding of the target entity, with the Full cascade condition, which contextualizes the target using topic routing and retrieved chunk-, sentence-, and entity-level evidence. As shown in Fig.~\ref{fig:same_entity_different_sense_disambiguation}, Raw matching cannot distinguish the two classes because identical entities receive identical static embeddings. In contrast, Full cascade achieves an AUROC of 0.774, with a 95\% confidence interval of $[0.691, 0.854]$. Its paired AUROC improvement over Raw is 0.274, with a 95\% bootstrap interval of $[0.191, 0.354]$. These results indicate that query-time cascaded contextualization can distinguish entities used with different meanings in the evaluated corpus.

\begin{figure}[H]
\centering
\includegraphics[width=0.82\linewidth]{figs/entity_contextualization_diagnostic.pdf}
\caption{Same-entity different-sense disambiguation diagnostic across document contexts. Different-sense entity pairs are treated as positive instances, and cosine distance is used as the disambiguation score. Points and horizontal bars show AUROC estimates and 95\% confidence intervals, respectively. Higher AUROC is better.}
\label{fig:same_entity_different_sense_disambiguation}
\end{figure}

\section{Query Representation Comparison for Chunk-level and Sentence-level Matching Scores in Sentence Reranking}
\label{app:query_scoring_comparison}

This appendix compares the query representations used to compute the chunk-level and sentence-level matching scores during sentence reranking. The main text uses the original representation of the expanded query for both scores. The other configurations remain unchanged.

We evaluate 50 first-stage questions per evaluation run from the mechanical-domain corpus using six metrics. Fig.~\ref{fig:query_scoring_comparison} shows that using the original representation of the expanded query for both chunk-level and sentence-level matching achieves the highest average score (86.48). The results indicate that retaining $\mathbf{q}^{(exp)}$ as the common semantic anchor is the most effective configuration.

\begin{center}
\includegraphics[width=0.72\linewidth]{figs/query_scoring_2x2_grouped_bar.pdf}
{\footnotesize
\captionof{figure}{Comparison of query representations for chunk-level and sentence-level matching scores on 50 first-stage mechanical-domain questions per evaluation run. Exp. denotes $\mathbf{q}^{(exp)}$ and Casc. denotes the corresponding cascaded query representation. Error bars denote the standard deviation across the five run-level mean scores, with each run containing 50 questions. Higher is better for all metrics.}
\label{fig:query_scoring_comparison}
}
\end{center}


\section{LLM Judge Scoring Prompt}
\label{app:llm_judge_scoring_prompt}

This appendix provides the scoring guidelines for the LLM judge. For each evaluation instance, the question, generated answer, and corresponding reference evidence are provided together with this rubric. The judge assigns one level and one numerical score within the corresponding range for each criterion.

\begin{lstlisting}
SCORING_PROMPT = """
You will evaluate the answers to the questions by using the relevant documents based on six criteria: **Comprehensiveness**, **Diversity**, **Empowerment**, **Logical**, **Readability**, and **Relevance**.

- **Comprehensiveness** -
Measure whether the answer comprehensively covers all key aspects of the question and whether there are omissions.
Level   | score range | description
Level 1 | 0-20   | The answer is extremely one-sided, leaving out key parts or important aspects of the question.
Level 2 | 20-40  | The answer has some content, but it misses many important aspects of the question and is not comprehensive enough.
Level 3 | 40-60  | The answer is more comprehensive, covering the main aspects of the question, but there are still some omissions.
Level 4 | 60-80  | The answer is comprehensive, covering most aspects of the question, with few omissions.
Level 5 | 80-100 | The answer is extremely comprehensive, covering all aspects of the question with no omissions.

- **Diversity** -
Measure the richness of the answer content, including background knowledge, extended information, case studies, etc.
Level   | score range | description
Level 1 | 0-20   | The answer is extremely sparse, providing only direct answers without additional information.
Level 2 | 20-40  | The answer provides a direct answer but contains only a small amount of relevant knowledge expansion.
Level 3 | 40-60  | In addition to the direct answers, the answer also provides some relevant background knowledge.
Level 4 | 60-80  | The answer is rich, providing more relevant background knowledge and supplementary information.
Level 5 | 80-100 | The answer provides a lot of relevant knowledge, expanded content and in-depth analysis.

- **Empowerment** -
Measure the credibility of the answer and whether it convinces the reader that it is correct.
Level   | score range | description
Level 1 | 0-20   | The answer lacks credibility, contains obvious errors or false information.
Level 2 | 20-40  | The answer has some credibility, but some information is not accurate.
Level 3 | 40-60  | The answer is credible and provides some supporting information.
Level 4 | 60-80  | The answer is highly credible, providing sufficient supporting information.
Level 5 | 80-100 | The answer is highly credible with authoritative supporting information.

- **Logical** -
Measure whether the answers are coherent, clear, and easy to understand.
Level   | score range | description
Level 1 | 0-20   | The answer is illogical, incoherent, and difficult to understand.
Level 2 | 20-40  | The answer has some logic, but is incoherent in parts.
Level 3 | 40-60  | The answer is logically clear and basically coherent.
Level 4 | 60-80  | The answer is logical, coherent, and easy to understand.
Level 5 | 80-100 | The answer is extremely logical, fluent and well-organized.

- **Readability** -
Measure whether the answer is well organized, clear in format, and easy to read.
Level   | score range | description
Level 1 | 0-20   | The format is confused and difficult to read.
Level 2 | 20-40  | There are some problems in the format.
Level 3 | 40-60  | The format is basically clear.
Level 4 | 60-80  | The format is clear and well organized.
Level 5 | 80-100 | The format is very clear with excellent reading experience.

- **Relevance** -
Measure whether the reasoning and answer are highly relevant and helpful to the question.
Level   | score range | description
Level 1 | 0-20   | The answer is entirely irrelevant, barely related to the question, or largely unhelpful.
Level 2 | 20-40  | The answer has limited relevance; much of the response is off-topic or unhelpful.
Level 3 | 40-60  | The answer is generally relevant, but includes distractions or less helpful parts.
Level 4 | 60-80  | The answer is mostly on point, with only minor digressions and overall useful content.
Level 5 | 80-100 | The answer is fully focused on the question, highly relevant, and highly helpful.

For each indicator, give a Level and a Score within the level's score range.
"""
\end{lstlisting}

\section{Detailed Experimental Setup}
\label{app:experimental_setup}

\subsection{Datasets}

We use eight document QA datasets: five machine-tool corpora (Mechanical, Hydraulics, Electrical, Maintenance, and System) and three general-domain corpora (Mix, Pathology, and Physics). The machine-tool corpora contain hierarchical technical descriptions and dense domain entities, whereas the three general-domain corpora test transfer under different terminology and document structures. Raw documents are deduplicated and converted into textual contexts. In each evaluation run, questions are generated from newly sampled corpus chunks, and their source contexts are retained as reference evidence.

Mechanical is the controlled benchmark for detailed analysis. In each evaluation run, it uses three 50-question stages: Stage~1 requires direct factual evidence, Stage~2 requires one intermediate evidence connection, and Stage~3 requires two connected evidence links. This progression tests retrieval from local support to compositional multi-hop evidence. Table~\ref{tab:dataset_statistics} reports corpus sizes and question counts per evaluation run.

\begin{table}[H]
\centering
\caption{Statistics of the experimental datasets.}
\label{tab:dataset_statistics}
\scriptsize
\renewcommand{\arraystretch}{0.92}
\begin{tabular*}{\linewidth}{@{\extracolsep{\fill}}llrr@{}}
\toprule
Dataset & Domain & Words & Questions / run \\
\midrule
Mechanical & Machine-tool & 3,161,292 & 50 / 50 / 50 \\
Hydraulics & Machine-tool & 2,145,104 & 50 \\
Electrical & Machine-tool & 2,189,922 & 50 \\
Maintenance & Machine-tool & 4,228,233 & 50 \\
System & Machine-tool & 3,291,093 & 50 \\
Mix & General-domain & 2,196,607 & 50 \\
Pathology & General-domain & 3,529,960 & 50 \\
Physics & General-domain & 2,369,660 & 50 \\
\bottomrule
\end{tabular*}
\end{table}

Word counts exclude whitespace after preprocessing and deduplication. The three Mechanical counts are the numbers of questions per evaluation run for the three stages. Within each run, all methods use the same documents, questions, reference evidence, generation model, and evaluation protocol. GPT-4o-mini is used for question generation and automatic answer evaluation.

\subsection{Baselines}

We compare with an LLM without retrieval, dense RAG \citep{Karpukhin2020DPR}, GraphRAG \citep{Edge2025GraphRAG}, LightRAG \citep{Guo2025LightRAG}, HippoRAG~2 \citep{Gutierrez2025HippoRAG2}, and Hyper-RAG \citep{Feng2026HyperRAG}. These baselines span unstructured retrieval, graph-based retrieval, memory-inspired graph retrieval, and hypergraph retrieval.

\subsection{Evaluation Metrics}

We use an LLM judge for open-ended knowledge-intensive QA \citep{Liu2023GEval,Zheng2023LLMJudge} and evaluate answers against the corresponding reference evidence. The six criteria are Comprehensiveness, Diversity, Empowerment, Logical, Readability, and Relevance, each scored from 0 to 100. All methods receive the same reference evidence, rubric, and judging prompt (Appendix~\ref{app:llm_judge_scoring_prompt}). We conduct five independent evaluation runs. In each run, we randomly sample different corpus chunks and generate a distinct 50-question evaluation set for each dataset and Mechanical stage. We report the mean $\pm$ standard deviation across the five run-level mean scores. The average score is the mean over the six criteria.

\section{Representative Mechanical Questions Across Evidence Depths}
\label{app:mechanical_question_examples}

Table~\ref{tab:mechanical_question_examples} presents representative questions from the three evidence-depth stages of the Mechanical dataset. Stage~1 requires direct factual support, Stage~2 introduces one intermediate evidence connection, and Stage~3 requires two connected evidence links. These examples illustrate how the evaluation progresses from retrieving local facts to assembling connected evidence for multi-hop answers.

\begin{table}[H]
\centering
\caption{Representative Mechanical-corpus questions across the three stages.}
\label{tab:mechanical_question_examples}
\setlength{\tabcolsep}{3pt}
\footnotesize
\renewcommand{\arraystretch}{1.10}
\begin{tabular}{@{}l>{\raggedright\arraybackslash}p{0.60\textwidth}>{\raggedright\arraybackslash}p{0.25\textwidth}@{}}
\toprule
Stage & \multicolumn{1}{c}{Representative question} & \multicolumn{1}{c}{Evidence requirement} \\
\midrule
Stage 1 & What factors influence the dynamic rigidity of spindle components in CNC machine tools? & Direct factual evidence. \\
Stage 2 & What are the key design requirements for CNC machine tool spindles to improve their precision and rigidity, and how do these requirements impact the machining accuracy of the workpiece? & An intermediate connection involving the spindle design requirements and their effect on machining accuracy. \\
Stage 3 & Explain how the design requirements for CNC machine tool spindles emphasize the need for high precision and rigidity, and specify the methods used to improve the rigidity of spindle components, particularly focusing on bearing configurations and preloading techniques. & Two connected evidence links involving spindle requirements, rigidity-improvement methods, and bearing design. \\
\bottomrule
\end{tabular}
\end{table}

\section{Expert Blind Evaluation}
\label{app:expert_blind_evaluation}

\begin{table}[H]
\centering
\caption{Expert blind evaluation on 60 paired Mechanical questions. Scores are question-level means over three experts on a 1--5 scale. $\Delta$ is AI-H\textsuperscript{2}RAG minus Hyper-RAG; CIs are 95\% paired-bootstrap CIs.}
\label{tab:mechanical_expert_evaluation}
\setlength{\tabcolsep}{5pt}
\small
\resizebox{0.94\linewidth}{!}{%
\begin{tabular}{lccccc}
\toprule
Metric & Hyper-RAG & AI-H\textsuperscript{2}RAG & $\Delta$ (Casc--Hyper) & 95\% CI & ICC(A,1) \\
\midrule
Correctness & 4.694 & \textbf{4.956} & +0.261 & [0.083, 0.489] & 0.757 \\
Completeness & 4.567 & \textbf{4.806} & +0.239 & [0.056, 0.456] & 0.770 \\
Clarity & 4.672 & \textbf{4.878} & +0.206 & [0.089, 0.333] & 0.765 \\
\bottomrule
\end{tabular}%
}
\end{table}

To complement the LLM judge, we conduct a blinded comparison of AI-H\textsuperscript{2}RAG and Hyper-RAG on 60 randomly sampled Mechanical questions. Three machine-tool experts independently receive the question, reference answer, and two anonymous system answers. The mapping to System~A or B and their presentation order are randomized for each expert and question, then revealed only after the forms are locked.

Experts score Correctness, Completeness, and Clarity on a 1--5 scale and select an overall preference. We average scores across experts for each question, estimate paired differences with 10,000 bootstrap resamples, and report 95\% confidence intervals. Inter-rater agreement is measured by two-way random-effects, absolute-agreement ICC(A,1). Table~\ref{tab:mechanical_expert_evaluation} shows higher scores for AI-H\textsuperscript{2}RAG in Correctness (4.956 vs. 4.694), Completeness (4.806 vs. 4.567), and Clarity (4.878 vs. 4.672). All paired-bootstrap intervals exclude zero. Of 180 paired judgments, experts preferred AI-H\textsuperscript{2}RAG in 51, Hyper-RAG in 26, and selected a tie in 103. This independent assessment is consistent with the automatic results.

\section{Key-evidence Coverage and Multi-hop Evidence-chain Retrieval}
\label{app:key_evidence_chain_retrieval}

\begin{table}[H]
\centering
\caption{Evidence coverage on the frozen Mechanical evaluation set. Values in the three recall columns are percentages, reported as AI-H\textsuperscript{2}RAG / Hyper-RAG. $\Delta$ is AI-H\textsuperscript{2}RAG minus Hyper-RAG in percentage points, with 95\% paired-bootstrap confidence intervals.}
\label{tab:rq6_evidence_chain_retrieval}
\setlength{\tabcolsep}{2pt}
\footnotesize
\renewcommand{\arraystretch}{1.12}
\resizebox{0.94\linewidth}{!}{%
\begin{tabular}{@{}lcccc@{}}
\toprule
Group & Sentence Recall & Bridge Recall & Full-chain Recall & $\Delta$ Full-chain Recall (95\% CI) \\
\midrule
Stage 1 & \textbf{95.0} / 50.0 & --- & \textbf{95.0} / 50.0 & \textbf{+45.0} [25.0, 65.0] \\
Stage 2 & \textbf{86.0} / 67.0 & \textbf{84.0} / 64.0 & \textbf{84.0} / 64.0 & \textbf{+20.0} [6.1, 34.7] \\
Stage 3 & \textbf{87.3} / 69.3 & \textbf{64.0} / 50.0 & \textbf{58.0} / 46.0 & \textbf{+12.0} [0.2, 24.0] \\
\bottomrule
\end{tabular}%
}
\renewcommand{\arraystretch}{1}

\raggedright\footnotesize Sentence Recall is the proportion of annotated gold source sentences visible in the final source evidence. Bridge Recall requires an annotated bridge entity and both gold source sentences linked by that entity to be visible. Full-chain Recall requires all annotated gold source sentences and bridge entities needed to form the complete evidence chain to be visible.
\end{table}

We finally test whether the answer-quality gain is accompanied by better retrieval of generator-visible evidence. Two annotators construct a frozen 120-question Mechanical subset for gold-evidence evaluation, comprising 20 questions from Stage~1, 50 from Stage~2, and 50 from Stage~3. Disagreements are adjudicated before evaluation. We map the native outputs of AI-H\textsuperscript{2}RAG and Hyper-RAG to the same original-document sentence identifiers and evaluate both under a shared 12,000-token source-evidence budget. Table~\ref{tab:rq6_evidence_chain_retrieval} shows that AI-H\textsuperscript{2}RAG outperforms Hyper-RAG on all applicable metrics across the three stages. Full-chain recall increases from 50.0\% to 95.0\% in Stage~1, from 64.0\% to 84.0\% in Stage~2, and from 46.0\% to 58.0\% in Stage~3. These improvements correspond to gains of 45.0, 20.0, and 12.0 percentage points, respectively, with all paired-bootstrap 95\% confidence intervals excluding zero. For Stages~2 and~3, bridge recall also improves by 20.0 and 14.0 percentage points, respectively. These results show that cascaded retrieval exposes both key source sentences and entity-mediated links more reliably to the generator.
\end{document}

